\subsection{MATRICES AND LINEAR MAPPINGS}

Let A be a ring and E a (right or left) A-module admitting a basis \((e_i)_{i\in I}\) . For every element \(x\in E\) , the matrix of \(x\) with respect to the basis \((e_i)\) , denoted by \(M(x)\) or \(\mathbf{x}\) (or sometimes simply \(x\) when no confusion can arise), is the column matrix consisting of the components \(x_i\) ( \(i\in I\) ) of \(x\) with respect to \((e_i)\) ( \(\S 1\) , no. 11); in calculations it will sometimes be convenient, in order to remember that the index \(i\) is a row index, to adjoin to it a column index taking only one value and write the matrix \(M(x)\) as \((x_{i0})\) .

We now consider two (left or right) A-modules E and F with bases \((e_{i})_{i\in I}\) and \((f_{k})_{k\in K}\) respectively; let \((f_{k}^{*})\) be the family of coordinate forms corresponding to \((f_{k})\) . For a linear mapping u of E into F, we shall define the matrix of u with respect to the bases \((e_{i})\) , \((f_{k})\) in each of the following cases:

\begin{enumerate}
\def\labelenumi{(\Alph{enumi})}
\setcounter{enumi}{3}
\item
  E and F are right A-modules, u is A-linear.
\item
  E and F are left A-modules, u is A-linear.
\end{enumerate}

In what follows, we shall attach the letter (D) (resp. (G)) to formulae applying to right (resp. left) modules.

DEFINITION 3. In each of the above two cases, the matrix of \(u\) with respect to the bases \((e_i), (f_k)\) is the matrix \(M(u) = (u_{ki})_{(k,i) \in K \times I}\) such that

\[
u _ {k i} = f _ {k} ^ {*} (u (e _ {i}))\tag{14}
\]

which is written respectively as

\[
u _ {k i} = \langle f _ {k} ^ {*}, u (e _ {i}) \rangle\tag{14 D}
\]

\[
u _ {k i} = \langle u (e _ {i}), f _ {k} ^ {*} \rangle .\tag{14 G}
\]

The column of \(M(u)\) of index i is therefore equal to \(M(u(e_{i}))\) .

Clearly if u, v are two linear mappings of E into F and \(M(u)\) , \(M(v)\) their matrices with respect to the same bases, then

\[
M (u + v) = M (u) + M (v)\tag{15}
\]

and

\[
M (\gamma u) = \gamma M (u)\tag{16}
\]

for every element \(\gamma\) of the centre \(\Gamma\) of A. In other words, once the bases \((e_i)\) , \((f_k)\) are fixed, the mapping \(u \mapsto M(u)\) is a \(\Gamma\) -module isomorphism of \(\operatorname{Hom}_{\mathbf{A}}(\mathbf{E}, \mathbf{F})\) onto a subset of the set \(\mathbf{A}^{\mathbf{K} \times \mathbf{I}}\) , equal to \(\mathbf{A}^{\mathbf{K} \times \mathbf{I}}\) if \(\mathbf{K}\) is finite.

PROPOSITION 2. Suppose I and K are finite. For every element \(x \in \mathbf{E}\) , the matrix \(M(u(x))\) with respect to the basis \((f_k)\) is given by the formula

\[
M (u (x)) = M (u). M (x)\tag{17 D}
\]

\[
{ } ^ { t } M ( u ( x ) ) = { } ^ { t } M ( x ) . { } ^ { t } M ( u ) .\tag{17 G}
\]

We verify for example (17 G). Let \(x = \sum_{i} x_{i0} e_i\) , \(u(x) = \sum_{k} y_{k0} f_k\) with \(x_{i0} \in \mathbf{A}\) , \(y_{k0} \in \mathbf{A}\) ; then \(u(x) = u\left(\sum_{i} x_{i0} e_i\right) = \sum_{i} x_{i0} u(e_i) = \sum_{i, k} x_{i0} u_{ki} f_k\) ; whence \(y_{k0} = \sum_{i} x_{i0} u_{ki}\) . In order to bring the two indices \(i\) along side one another, we consider the transpose matrices \(^t M(x) = (x'_{0i})\) , where \(x'_{0i} = x_{i0}\) and

\[
{ } ^ { t } M ( u ) = \left( u _ { i k } ^ { \prime } \right) ,
\]

where \(u_{ik}' = u_{ki}\) ; then \(y_{k0} = \sum_{i} x_{0i}' u_{ik}'\) and the right hand side is the element of index \(k\) of the matrix with one row \(^t M(x).^t M(u)\) , whence (17 G).

When A is commutative, (17 G) reduces to (17 D) by formula (4) of no. 2.

COROLLARY. Let E, F, G be three right (resp. left) modules over a ring A, \((e_{i})_{i\in I}\) , \((f_{k})_{k\in K}\) , \((g_{l})_{l\in L}\) respective finite bases of E, F, G, \(u:E\to F\) , \(v:F\to G\) two linear mappings, \(M(u)\) the matrix of u relative to the bases \((e_{i})\) , \((f_{k})\) , \(M(v)\) the matrix of v relative to the bases \((f_{k})\) , \((g_{l})\) and \(M(v\circ u)\) the matrix of \(v\circ u\) relative to the bases \((e_{i})\) , \((g_{l})\) ; then

\[
M (v \circ u) = M (v) M (u)\tag{18 D}
\]

\[
{ } ^ { t } M ( v \circ u ) = { } ^ { t } M ( u ) { } ^ { t } M ( v ) .\tag{18 G}
\]

We prove for example (18 G). For all \(x \in \mathbf{E}\) , by (17 G):

\(^{t}M(x).^{t}M(v\circ u)=^{t}M(v(u(x)))=^{t}M(u(x)).^{t}M(v)=^{t}M(x).^{t}M(u).^{t}M(v)\) by associativity; the corollary then follows from no. 3, Proposition 1 since the matrix \(^{t}M(x)\) with one row is arbitrary.

Remark (1). Formula (17 D) can be considered as a special case of (18 D). For there corresponds canonically to every \(x \in \mathbf{E}\) the linear mapping \(\theta_x: \mathbf{A}_d \to \mathbf{E}\) mapping every \(\alpha \in \mathbf{A}\) to \(x\alpha\) (§ 2, no. 1). It is immediate that the matrix \(M(\theta_x)\) with respect to the basis 1 of \(\mathbf{A}_d\) and the basis \((e_i)\) of \(\mathbf{E}\) is just the matrix \(M(x)\) ; similarly \(M(\theta_{u(x)}) = M(u(x))\) and formula (17 D) can therefore be considered as a translation of the relation

\[
\theta_ {u (x)} = u \circ \theta_ {x}.
\]

PROPOSITION 3. Let E, F be two right (resp. left) A-modules and \((e_i)_{i\in I}, (f_k)_{k\in K}\) finite bases of E and F respectively. For every linear mapping \(u\) of E into F, let \(M(u)\) be the matrix of \(u\) with respect to the bases \((e_i)\) and \((f_k)\) . Then the matrix of \(^t u: F^* \to E^*\) with respect to the dual bases \((f_k^*)\) and \((e_i^*)\) is equal to \(^t M(u)\) .

E is canonically identified with its bidual \(E^{**}\) and \((e_{i})\) with the dual basis of \((e_{i}^{*})\) ; then (supposing for example that E and F are right modules)

\[
\langle^ {t} u (f _ {k} ^ {*}), e _ {i} \rangle = \langle f _ {k} ^ {*}, u (e _ {i}) \rangle ,
\]

whence the proposition.

Remarks. (2) Let E and F be two left A-modules with bases \((e_{i})_{i\in\mathbf{I}}\) and \((f_{k})_{k\in\mathbf{K}}\) respectively. For every A-linear mapping \(u:E\to F\) , by (14 G), \(u(e_{i})=\sum_{k}u_{ki}f_{k}\) ; these relations can also be interpreted by saying that the column matrix \((u(e_{i})_{i\in\mathbf{I}})\) with elements in F is equal to the product \({}^{t}\mathbf{M}(u).(f_{k})\) , where \((f_{k})_{k\in\mathbf{K}}\) is considered as a column matrix with elements in F and the product is calculated for the mapping \(A\times F\to F\) defining the law of action on the A-module F (no. 2).

\begin{enumerate}
\def\labelenumi{(\arabic{enumi})}
\setcounter{enumi}{2}
\tightlist
\item
  Let A, B be two commutative rings and \(\sigma: A \to B\) a ring homomorphism. In the notation of Proposition 3, \((e_i \otimes 1)\) and \((f_k \otimes 1)\) are respective bases of \(E_{(B)} = E \otimes_A B\) and \(F_{(B)} = F \otimes_A B\) ( \(\S 5\) , no. 1, Proposition 4); moreover, if \((e_i^*)\) and \((f_k^*)\) are respectively the dual bases of \((e_i)\) and \((f_k)\) , then \((e_i^* \otimes 1)\) and \((f_k^* \otimes 1)\) are respectively the dual bases of \((e_i \otimes 1)\) and \((f_k \otimes 1)\) ( \(\S 5\) , no. 4). For every A-linear mapping \(u: E \to F\) , let \(M(u)\) and \(M(u \otimes 1)\) be the matrix of \(u\) with respect to \((e_i)\) and \((f_k)\) and the matrix of the B-linear mapping \(u \otimes 1\) with respect to \((e_i \otimes 1)\) and \((f_k \otimes 1)\) . It follows from \(\S 5\) , no. 4, formula (20) that
\end{enumerate}

\[
M (u \otimes 1) = \sigma (M (u)).
\]

Consider a system of a finite number of right scalar linear equations in a finite number of unknowns

\[
\sum_ {i \in I} a _ {k i} x _ {i} = b _ {k} \quad (k \in K)\tag{19}
\]

with \(a_{ki}\) , \(x_{i}\) , \(b_{k}\) in A.

Let \((e_i)_{i\in I}, (f_k)_{k\in K}\) be the canonical bases of \(\mathbf{E} = \mathbf{A}_d^{\mathrm{I}}\) and \(\mathbf{F} = \mathbf{A}_d^{\mathrm{K}}\) ; the system (19) is equivalent to the equation \(u(x) = b\) , where \(x = \sum_{i} e_i x_i\) , \(b = \sum_{k} f_k b_k\) and \(u: \mathbf{E} \to \mathbf{F}\) is the linear mapping such that the matrix \(M(u)\) with respect to the bases \((e_i)\) and \((f_k)\) are equal to \(\mathbf{A} = (a_{ki})_{(k,i)\in K\times I}\) . This matrix is called the matrix of the system of linear equations (19). Recall (§ 2, no. 8, Remarks 2 and 3), that, writing \(c_i = \sum_{k} f_k a_{ki}\) , the system (19) is equivalent to the unique vector equation

\[
\sum_ {i} c _ {i} x _ {i} = b,\tag{20}
\]

and as \(c_{i}\) is the column of index i in the matrix A, we see that to say that the system (19) admits a solution amounts to saying that the matrix \(b = (b_{k0})\) with one column is a linear combination of the columns of the matrix A.

We leave to the reader the task of formulating the analogous definitions and remarks for systems of left linear equations.

